\documentclass{article}

% -------------------------------------------------
% CSE 579: Knowledge Representation and Reasoning
% Course Project Report Template (NeurIPS style)
% -------------------------------------------------
% Main report length: 2--3 pages
% Contribution table: required at the end, NOT included in page limit
% Use this template with the provided neurips_2026.sty file.
%
% Note to students:
% - Keep the report concise and focused.
% - Your report should clearly show:
%   (1) the problem,
%   (2) the technical system design,
%   (3) the empirical evaluation and analysis.
% - The contribution table should appear after the main report.

\usepackage[main,final]{neurips_2026}

\usepackage[utf8]{inputenc}
\usepackage[T1]{fontenc}
\usepackage{hyperref}
\usepackage{url}
\usepackage{booktabs}
\usepackage{amsfonts}
\usepackage{nicefrac}
\usepackage{microtype}
\usepackage{xcolor}
\usepackage{graphicx}
\usepackage{amsmath,amssymb}
\usepackage{enumitem}
\usepackage{multirow}
\usepackage{array}

\setlist[itemize]{leftmargin=1.3em, topsep=2pt, itemsep=1pt, parsep=0pt}
\setlist[enumerate]{leftmargin=1.5em, topsep=2pt, itemsep=1pt, parsep=0pt}

\title{CSE 579 Course Project Report\\[0.2em]
\large Project Title}

% Replace names as needed.
% Keeping all students in a single centered block avoids affiliation alignment issues.
\author{%
Student Name 1 \quad Student Name 2 \quad Student Name 3 \\
Arizona State University \\
CSE 579: Knowledge Representation and Reasoning
}

\begin{document}

\maketitle

\begin{abstract}
Briefly summarize the problem, your proposed system, how you evaluated it, and the
main takeaway in 4--6 sentences.
\end{abstract}

\section{Problem and Motivation}
State the problem clearly and explain why it matters. Mention which course project
direction this report relates to, or note that it is an approved custom project.
Define the task, inputs/outputs, and the main challenge.

\textbf{Suggested content.}
\begin{itemize}
    \item Problem statement
    \item Why the problem matters
    \item Relation to knowledge representation, reasoning, or multimodal understanding
    \item Main research question or hypothesis
\end{itemize}

\section{Method / System Design}
Explain the technical system you built. Focus on the main components, how they
interact, and what design choices you made. The report should make clear that the
project goes beyond prompting alone and includes an actual system design.

\textbf{Suggested content.}
\begin{itemize}
    \item System overview
    \item Knowledge representation used (e.g., graph, memory, retrieved text, uncertainty, tools)
    \item Reasoning, retrieval, verification, or routing mechanism
    \item Baseline or comparison method(s)
\end{itemize}

You may include one small figure or pipeline diagram if it helps.

\section{Experimental Setup}
Describe how you evaluated the system.

\textbf{Suggested content.}
\begin{itemize}
    \item Dataset, benchmark, or collected examples
    \item Evaluation metrics
    \item Baselines
    \item Important implementation details
\end{itemize}

\section{Results and Analysis}
Present the main findings. A compact table or figure is encouraged if it helps.
Do not only report numbers---also explain what the results mean.

\textbf{Suggested content.}
\begin{itemize}
    \item Main quantitative or qualitative results
    \item Comparison to baseline(s)
    \item Failure cases or limitations
    \item Key insight learned from the experiments
\end{itemize}

\section{Conclusion}
Summarize the main takeaway in a short paragraph. Briefly mention one limitation or
next step if relevant.

\section*{References}
Use a consistent citation style and cite all datasets, models, baselines, and prior work.
Keep references concise.

Please note that references are NOT included in the page limit.

% Example reference format:
% [1] Author names. Paper title. Venue, year.

\newpage
\section*{Contribution Table}
This is a mandatory section and not included in page limit.
\begin{table}[h]
    \centering
    \small
    \begin{tabular}{p{0.18\linewidth} p{0.57\linewidth} p{0.15\linewidth}}
        \toprule
        \textbf{Member} & \textbf{Contribution} & \textbf{Share / Role} \\
        \midrule
        Student 1 & Example: problem formulation, system design, implementation, writing & 35\% \\
        Student 2 & Example: data preparation, baseline implementation, experiments & 35\% \\
        Student 3 & Example: evaluation, analysis, presentation, editing & 30\% \\
        \bottomrule
    \end{tabular}
\end{table}

\end{document}
